\addtocontents{toc}{\protect\Needspace{3\baselineskip}}
\originalchapter{距离、吸收与观测推断}\label{ch:distance}
\originalsection{三种距离}
由径向光线得共动径向距离 $\chi(z)=cH_0^{-1}\int_0^z\dd z'/E(z')$. 曲空间横向距离为 $D_M=a_0S_k(\chi)$, 其中 $S_k$ 分别为球面、平直或双曲几何的正弦、恒等、双曲正弦形式, 并保留相应曲率半径. 以下数值题的默认平直情形有 $D_M=\chi$.

发射时大小为 $l$ 的物体张角 $\theta=l/D_A$, $D_A=D_M/(1+z)$. 每个光子能量下降 $1+z$, 到达间隔再增加 $1+z$, 所以 $F=L/[4\pi D_M^2(1+z)^2]$. 定义 $D_L$ 使 $F=L/(4\pi D_L^2)$, 得
\begin{equation}\label{eq:distances}
D_L=(1+z)D_M=(1+z)^2D_A.
\end{equation}
此关系需要光子数守恒和正常的度规光传播; 吸收、透镜选择效应或仪器选择会影响实际测量. 高红移物体的角大小不能用今日实际距离直接除, 宇宙学距离也不能统一称为一个「距离」而省掉定义.

\originalsection{超新星与参数退化}
Ia 型超新星通过光变宽度和颜色修正构造距离模数. 典型参数化为 $\mu_i=m_i^*-M+\alpha(s_i-1)-\beta c_i$, 与 $\mu_{\rm th}=5\log_{10}[D_L/(\mathrm{Mpc})]+25$ 比较. $M$ 与 $H_0$ 存在归一化退化, 不能只凭超新星相对距离各自精确确定两者. 若 $\sigma_i^2$ 依赖 $\alpha,\beta$, 简单只最小化 $\sum(\mu_i-\mu_{\rm th})^2/\sigma_i^2$ 会偏向放大分母; 需要完整统计模型或题面指定的迭代权重策略.

低红移展开为 $D_L=(cz/H_0)[1+(1-q_0)z/2+O(z^2)]$, $q_0=\Omega_{m0}/2-\Omega_{\Lambda0}$ 对尘埃与真空成立. 很低红移只确定第一项, 较高红移才能比较减速或加速历史. 星系间尘埃使超新星更暗, 可模拟更大距离; 颜色、红外、背景辐射和不同距离探针的比较为这种退化提供检验. 本卷不捏造没有输入数据的最佳拟合.

\originalsection{吸收统计与 Lyman--alpha 森林}
沿实际路径 $\dd l=c\dd z/[(1+z)H(z)]$, 共动密度恒定 $n_c$ 的吸收体有实际密度 $n_c(1+z)^3$. 对固定实际截面 $\sigma$, 遭遇概率给
\begin{equation}\label{eq:absorber}
\frac{\dd N}{\dd z}=n_c\sigma\frac{c(1+z)^2}{H(z)}.
\end{equation}
若指定固定共动截面, 则实际截面还按 $(1+z)^{-2}$ 变, 红移趋势因此改变. 题面中的截面约定必须先确定.

氢的 Lyman--alpha 共振发生在约 $1216$ 埃. 远处类星体光穿过多个中性氢区域, 不同位置的共振因红移落在不同观测波长, 形成森林. 在均匀 Hubble 流的窄线近似下, 由 $\tau_\nu=\int n_{\rm HI}\sigma_\nu\dd l$ 与频率沿程扫描的速率可得 Gunn--Peterson 光学深度
$\tau_{\rm GP}=\pi e^2f_\alpha\lambda_\alpha n_{\rm HI}/[m_ecH(z)]$, 此式使用高斯电磁单位. 即使中性分数很小也可能强烈吸收. 光电离平衡 $n_{\rm HI}\Gamma\simeq\alpha(T)n_en_p$ 使森林同时依赖密度、温度和电离背景, 因而吸收不是单一的密度计; 速度与热展宽也改变谱线.
